\section{Implementation}\label{sec:impl}
The web application is a single HTML file of about 4500 lines with 241 JavaScript functions; the command-line program is a single Python file of about 4000 lines with 95 functions. Plotly.js 2.27 renders all interactive figures, including 3D isosurfaces; a plot editor stores per-curve colours, names, line styles and visibility across tabs. Figures are exported at publication resolution as PNG, JPEG, WebP, SVG or JPEG~XL, individually or all at once in a ZIP archive; JPEG~XL is encoded in the browser by the reference libjxl compiled to WebAssembly (lossless mode verified to be pixel-identical to PNG at about two thirds of its size) and on the command line through imagecodecs, where \code{--format png,jxl,svg,pdf} writes every figure in all requested formats. Other exports include tab-separated tables with a header line (recognised by Origin as long names), CHGCAR files in VASP~5 format with element names (as $\rho V$, readable by VESTA and \code{bader}), XSF files and ZIP archives with one folder per structure.

\paragraph{Internationalisation.} The interface, help and tutorial are available in six languages (Portuguese, English, Spanish, French, Italian and Latin). A small engine translates rendered text nodes and attributes through a dictionary of about 730 entries plus regular-expression rules for strings containing numbers (e.g.\ ``1.3942 indirect''); it stores the original text of each node, observes the document with a \code{MutationObserver} so that content produced later by the analyses is also translated, and ignores elements marked \code{translate="no"}. Numerical results are never altered. Long documents (the tutorial) have one page per language, generated from translated bodies, with \code{hreflang} metadata, and PDF versions are produced from the same pages with pandoc and XeLaTeX.

\paragraph{Packaging and distribution.} The Python package (\code{pip install brumare}) provides the commands \code{brumare} (command line), \code{brumare-gui} (desktop application) and \code{brumare-lote} (batch analysis with a minimal window). Continuous integration builds the Windows executables and installer for every tagged release and publishes the package to PyPI through trusted publishing. The software is released under the GPL-3.0-or-later \cite{gplv3}; bundled third-party components (Plotly.js and JSZip under MIT, MathJax under Apache~2.0, fonts under the SIL Open Font License) are compatible with it.

\paragraph{Testing.} An automated test suite (80 \code{pytest} tests on real outputs of the five codes, covering band gaps, effective masses, magnetic moments, the band-gap correction, lossless JPEG~XL export, Bader charge conservation, citation output and the translated command-line help, plus the cross-validation of Table~\ref{tab:bench}) runs under continuous integration on Python 3.9 and 3.12.
